\documentclass[11pt]{article}
\usepackage{notes} % place notes.sty in the same folder
\definecolor{AccN}{HTML}{A8861A}
\fancyhead[L]{\small\color{NavyN}\nouppercase{\leftmark}}\fancyhead[R]{\small\color{GreyN}Instructor Manual --- Lecture 10}
\usepackage{longtable,multicol,needspace}
\setlength{\emergencystretch}{2em}
\newtcolorbox{instr}[1]{nb,colback=AccN!6,colframe=AccN,title={#1}}
\newtcolorbox{prob}{nb,colback=BluN!5,colframe=BluN,title={Problem statement}}
\newtcolorbox{soln}{nb,colback=GrnN!5,colframe=GrnN,title={Model solution and marking points}}
\newtcolorbox{sayit}{nb,colback=white,colframe=NavyN,title={Read aloud to students}}
\newtcolorbox{live}[1]{nb,colback=AccN!4,colframe=AccN,title={#1}}
\begin{document}
\begin{titlepage}\centering\vspace*{2.2cm}
{\color{AccN}\large\bfseries LECTURE 10\par}\vspace{0.6cm}
{\color{NavyN}\Huge\bfseries Project Presentations \&\\Q\&A Panel\par}\vspace{0.6cm}
{\color{AccN}\rule{0.7\textwidth}{1.5pt}}\par\vspace{0.6cm}
{\LARGE Instructor Manual\par}\vspace{0.4cm}
{\large\begin{minipage}{0.85\textwidth}\centering Applying AI \& ML to a defence problem of your choosing\\Structure \textbullet\ rubric \textbullet\ project catalogue \textbullet\ Q\&A panel guide \textbullet\ course conclusion\end{minipage}\par}\vspace{2cm}
{\large AK DHAMIJA, Researcher \& Instructor, AI \& Machine Learning\par}\vspace{0.2cm}
{\normalsize \href{mailto:dhamija.ak@gmail.com}{dhamija.ak@gmail.com} \quad\textbullet\quad 9811453163\par}\vspace{0.3cm}
{Unit 4: Simulation Strategies in Warfare --- AI \& Machine Learning in Warfare\par}\vfill
\begin{tcolorbox}[colback=NavyN!4,colframe=NavyN,boxrule=0.6pt,arc=2pt,width=0.9\textwidth]\small
\textbf{This manual contains:} the 45-minute session plan with timings; a slide-by-slide teaching guide for all 11 slides; key messages, facilitation tips, anticipated questions with model answers, misconceptions and in-class activities; live demonstrations with the interactive studio; consistency notes on the source material; the complete quiz answer key with rationale and scoring rules; and the assignment with marking scheme and full model solutions.
\end{tcolorbox}
\end{titlepage}
\tableofcontents\newpage


\needspace{12\baselineskip}\section{Lecture at a Glance}

\begin{center}\small\renewcommand{\arraystretch}{1.25}\begin{tabularx}{\textwidth}{@{}>{\bfseries}p{3.3cm}>{\raggedright\arraybackslash}X@{}}\toprule

\textcolor{NavyN}{Lecture}&10 --- Project Presentations \& Q\&A Panel\\
\textcolor{NavyN}{Duration}&45 minutes: 5 min opening Q\&A, 30 min teaching, 10 min auto-quiz\\
\textcolor{NavyN}{Slides}&11 (Beamer PDF, 16:9), 8 sections\\
\textcolor{NavyN}{Prerequisite}&Lecture~9; students should have attempted its quiz and assignment.\\
\textcolor{NavyN}{Student materials}&\texttt{Lecture10\_Slides\_Beamer.pdf}; \texttt{Lecture10\_Reference\_Text.pdf}; \texttt{Lecture10\_Quiz.html}; \texttt{Lecture10\_Answers\_Assignment.html} (release after the quiz)\\
\textcolor{NavyN}{Assessment}&Auto-quiz: 15 MCQs, 40~s each, $+1/-0.25$. Assignment: 3 problems, 100 marks.\\\bottomrule\end{tabularx}\end{center}

\subsection*{Overview}Lecture~10 is the project presentation and Q\&A panel session. It explains how the session runs, the ten-minute, six-section presentation framework, the assessment rubric (problem definition and context 20\%, algorithm justification 25\%, results and validation 25\%, operational applicability 20\%, critical evaluation and ethics 10\%; peer evaluation contributes 20\% of the project grade), a catalogue of eight suggested topics, how the panel questions (judgement, not recall), the STARE framework for tough questions, common pitfalls, and lessons from other disciplines on presenting evidence responsibly. It closes the course with the arc of ten lectures and a reading list.

\subsection*{Learning outcomes}By the end of the lecture students should be able to:\begin{enumerate}\item Present an AI/ML defence project as a coherent, evidence-based argument.\item Report results with baselines, appropriate metrics and uncertainty.\item Handle panel questions honestly using STARE.\item State limitations and deployment conditions, including autonomy level and IHL considerations.\end{enumerate}

\subsection*{Syllabus mapping}\begin{itemize}\item Application of all Unit~4 syllabus items in a student project\end{itemize}

\needspace{12\baselineskip}\section{Session Plan (45 Minutes)}

\begin{center}\begin{tikzpicture}[x=0.3444cm,y=1cm]
\fill[AmbN!40](0,0)rectangle(5,1);\node[font=\scriptsize\bfseries,align=center]at(2.5,0.5){Q\&A\\recap};
\fill[AccN!25](5.00,0)rectangle(8.00,1);\draw[white,thick](5.00,0)--(5.00,1);
\node[font=\tiny,rotate=90,anchor=center]at(6.50,0.5){\parbox{0.9cm}{\centering \S1}};
\fill[NavyN!18](8.00,0)rectangle(11.00,1);\draw[white,thick](8.00,0)--(8.00,1);
\node[font=\tiny,rotate=90,anchor=center]at(9.50,0.5){\parbox{0.9cm}{\centering \S2}};
\fill[AccN!25](11.00,0)rectangle(14.00,1);\draw[white,thick](11.00,0)--(11.00,1);
\node[font=\tiny,rotate=90,anchor=center]at(12.50,0.5){\parbox{0.9cm}{\centering \S3}};
\fill[NavyN!18](14.00,0)rectangle(17.00,1);\draw[white,thick](14.00,0)--(14.00,1);
\node[font=\tiny,rotate=90,anchor=center]at(15.50,0.5){\parbox{0.9cm}{\centering \S4}};
\fill[AccN!25](17.00,0)rectangle(23.00,1);\draw[white,thick](17.00,0)--(17.00,1);
\node[font=\tiny,rotate=90,anchor=center]at(20.00,0.5){\parbox{0.9cm}{\centering \S5}};
\fill[NavyN!18](23.00,0)rectangle(26.00,1);\draw[white,thick](23.00,0)--(23.00,1);
\node[font=\tiny,rotate=90,anchor=center]at(24.50,0.5){\parbox{0.9cm}{\centering \S6}};
\fill[AccN!25](26.00,0)rectangle(29.00,1);\draw[white,thick](26.00,0)--(26.00,1);
\node[font=\tiny,rotate=90,anchor=center]at(27.50,0.5){\parbox{0.9cm}{\centering \S7}};
\fill[NavyN!18](29.00,0)rectangle(35.00,1);\draw[white,thick](29.00,0)--(29.00,1);
\node[font=\tiny,rotate=90,anchor=center]at(32.00,0.5){\parbox{0.9cm}{\centering \S8}};
\fill[RedN!30](35,0)rectangle(45,1);\node[font=\scriptsize\bfseries,align=center]at(40,0.5){Auto-quiz\\15 Q $\times$ 40 s};
\draw[NavyN](0,0)rectangle(45,1);
\foreach \tt in {0,5,10,15,20,25,30,35,40,45}{\draw[NavyN](\tt,0)--(\tt,-0.12);\node[font=\tiny,below]at(\tt,-0.12){\tt};}
\node[font=\scriptsize,below]at(22.5,-0.45){minutes};
\end{tikzpicture}\end{center}

\begin{center}\small\renewcommand{\arraystretch}{1.2}\begin{tabularx}{\textwidth}{@{}l l X r@{}}\toprule\textbf{Clock} & \textbf{Slides} & \textbf{Segment} & \textbf{Min}\\\midrule

0:00--5:00 & --- & \textbf{Opening Q\&A and recap} (see prompts below) & 5\\

5:00--8:00 & 2 & 1. Session (1 slide) & 3\\

8:00--11:00 & 3 & 2. Presentation (1 slide) & 3\\

11:00--14:00 & 4 & 3. Rubric (1 slide) & 3\\

14:00--17:00 & 5 & 4. Topics (1 slide) & 3\\

17:00--23:00 & 6--7 & 5. Q\&A Panel (2 slides) & 6\\

23:00--26:00 & 8 & 6. Pitfalls (1 slide) & 3\\

26:00--29:00 & 9 & 7. Across Disciplines (1 slide) & 3\\

29:00--35:00 & 10--11 & 8. Conclusion (2 slides) & 6\\

35:00--45:00 & --- & \textbf{Auto-quiz} (instructions 1 min, 15 questions $\times$ 40 s = 10 min of answering; collect scores) & 10\\\bottomrule\end{tabularx}\end{center}

Slide 1 is the title slide and is shown while students settle. The teaching segment covers 10 content slides in 30 minutes, an average pace of 0.3 slides per minute. 

The pace allows time for the in-class activities in Section~\ref{sec:act}.

\begin{instr}{Opening Q\&A prompts (first 5 minutes)}\begin{enumerate}\item From Lecture~9: what are the four questions of the algorithm-selection decision tree?\item Name two exam traps from the misconceptions slides.\item (Logistics) Confirm presentation order and timing.\end{enumerate}Take two or three answers per prompt; correct misconceptions briefly and note any topic that needs revisiting.\end{instr}

\needspace{12\baselineskip}\section{Slide-by-Slide Teaching Guide}

Each subsection gives the slide range, the suggested time, what to emphasise, and the slide titles in order. Numbers refer to pages of the Beamer PDF.

\needspace{8\baselineskip}\subsection{1. Session \normalfont\small\color{GreyN}(slide 2, about 3 min)}

\textbf{Teaching note.} The deck describes a three-hour capstone: Hour~1, groups of 4--5 give a 10-minute talk plus 2 minutes of panel Q\&A (slots by lot); Hour~2, open-floor panel deep-dive; Hour~3, feedback, industry perspectives and reading. In a 45-minute slot, run three groups and move the open floor to a later session.

\begin{center}\small\renewcommand{\arraystretch}{1.1}\begin{tabularx}{\textwidth}{@{}r X@{}}2 & Lecture 10 --- How the Session Runs\\\end{tabularx}\end{center}

\needspace{8\baselineskip}\subsection{2. Presentation \normalfont\small\color{GreyN}(slide 3, about 3 min)}

\textbf{Teaching note.} The six-section framework: problem, data, method, results, deployment, limitations and ethics.

\begin{center}\small\renewcommand{\arraystretch}{1.1}\begin{tabularx}{\textwidth}{@{}r X@{}}3 & The Ten-Minute Framework --- Six Sections, One Story\\\end{tabularx}\end{center}

\needspace{8\baselineskip}\subsection{3. Rubric \normalfont\small\color{GreyN}(slide 4, about 3 min)}

\textbf{Teaching note.} Walk through the rubric bands (distinction, merit, pass, fail) and weights; explain peer evaluation.

\begin{center}\small\renewcommand{\arraystretch}{1.1}\begin{tabularx}{\textwidth}{@{}r X@{}}4 & Assessment Rubric --- Five Criteria, Each Marked Out of 20\\\end{tabularx}\end{center}

\needspace{8\baselineskip}\subsection{4. Topics \normalfont\small\color{GreyN}(slide 5, about 3 min)}

\textbf{Teaching note.} Eight suggested topics, each mapped to the lecture that supplies the method.

\begin{center}\small\renewcommand{\arraystretch}{1.1}\begin{tabularx}{\textwidth}{@{}r X@{}}5 & Project Topic Catalogue --- Eight Suggested Themes\\\end{tabularx}\end{center}

\needspace{8\baselineskip}\subsection{5. Q\&A Panel \normalfont\small\color{GreyN}(slides 6--7, about 6 min)}

\textbf{Teaching note.} What the panel asks and why; the STARE framework (State, Tie, Acknowledge, Reason, End).

\begin{center}\small\renewcommand{\arraystretch}{1.1}\begin{tabularx}{\textwidth}{@{}r X@{}}6 & The Q\&A Panel --- It Tests Judgement, Not Algorithm Recall\\
7 & Handling Tough Questions --- the STARE Framework\\\end{tabularx}\end{center}

\needspace{8\baselineskip}\subsection{6. Pitfalls \normalfont\small\color{GreyN}(slide 8, about 3 min)}

\textbf{Teaching note.} Common pitfalls: no baseline, accuracy on imbalanced data, random splits on sequential data, ``ready for deployment'' from simulation.

\begin{center}\small\renewcommand{\arraystretch}{1.1}\begin{tabularx}{\textwidth}{@{}r X@{}}8 & Common Project Pitfalls --- Avoid These in Your Presentation\\\end{tabularx}\end{center}

\needspace{8\baselineskip}\subsection{7. Across Disciplines \normalfont\small\color{GreyN}(slide 9, about 3 min)}

\textbf{Teaching note.} How medicine (trial reporting), aviation (accident reports) and psychology (replication) present evidence responsibly.

\begin{center}\small\renewcommand{\arraystretch}{1.1}\begin{tabularx}{\textwidth}{@{}r X@{}}9 & Presenting Evidence Responsibly --- Lessons From Other Disciplines\\\end{tabularx}\end{center}

\needspace{8\baselineskip}\subsection{8. Conclusion \normalfont\small\color{GreyN}(slides 10--11, about 6 min)}

\textbf{Teaching note.} The arc of ten lectures and recommended reading.

\begin{center}\small\renewcommand{\arraystretch}{1.1}\begin{tabularx}{\textwidth}{@{}r X@{}}10 & Course Conclusion --- The Arc of Ten Lectures\\
11 & Recommended Reading --- Going Further\\\end{tabularx}\end{center}

\needspace{12\baselineskip}\section{Key Messages}Students should leave the room able to state these in their own words:\begin{keypoint}\begin{enumerate}\item A project is an argument: problem, evidence, limits.\item Baselines, the right metric and uncertainty are non-negotiable.\item Honest ``I don't know, but here is how I would find out'' beats bluffing.\item Deployment means autonomy level, latency, interface, ROE/IHL compliance and monitoring.\end{enumerate}\end{keypoint}

\needspace{12\baselineskip}\section{Facilitation Tips}\begin{itemize}\item Keep strict time (10 + 2 minutes per group); a visible countdown helps. Invite 1--2 external defence-AI practitioners to the panel if possible.\item Panel: ask one question each on justification, validation and deployment/ethics; avoid recall questions.\item Collect peer scores on paper or a form immediately after each talk to avoid drift.\item If the session has more groups than fit, run the quiz in the first slot of the next session instead.\end{itemize}

\needspace{12\baselineskip}\section{Anticipated Questions and Model Answers}

\begin{question}{anticipated question}\textbf{Q:} Can we use a public dataset unrelated to defence?\par\smallskip\textbf{A:} Yes, if you argue the analogy explicitly and discuss what would differ with operational data.\end{question}

\begin{question}{anticipated question}\textbf{Q:} What if our model performs badly?\par\smallskip\textbf{A:} A well-diagnosed negative result with honest analysis can earn a distinction; unjustified claims cannot.\end{question}

\needspace{12\baselineskip}\section{Common Misconceptions}\begin{center}\small\renewcommand{\arraystretch}{1.35}\begin{tabularx}{\textwidth}{@{}>{\raggedright\arraybackslash}X >{\raggedright\arraybackslash}X@{}}\toprule\textbf{Misconception} & \textbf{Correction}\\\midrule``More slides of results make a better talk.'' & One clear results slide with a baseline and uncertainty is better.\\
``Ethics is a closing slide.'' & It is assessed and should shape design choices throughout.\\\bottomrule\end{tabularx}\end{center}

\needspace{12\baselineskip}\section{In-Class Activities and Discussion Prompts}\label{sec:act}\begin{itemize}\item \textbf{Mock panel (if time before presentations):} two volunteers answer a tough question using STARE; the class critiques.\end{itemize}

\needspace{12\baselineskip}\section{Run It Live: Studios as Project Resources}\label{sec:live}
Every lecture from 1 to 9 has an interactive studio in the course folder. Each is a single, self-contained HTML file whose code students may read, adapt and extend. Encourage project teams to use them in two ways: as a \textbf{live demonstration} inside the presentation (two minutes, one setting changed, one prediction tested), and as a \textbf{starting point} for a prototype, for example adding a new dataset to the clustering lab or a new map to the gridworld.
\begin{center}\small\renewcommand{\arraystretch}{1.2}\begin{tabularx}{\textwidth}{@{}c l X@{}}\toprule\textbf{L} & \textbf{File} & \textbf{Useful for projects on}\\\midrule
1 & \texttt{Lecture1\_Data\_Can\_Lie\_Studio.html} & evaluation pitfalls, base rates, multiple testing\\
2 & \texttt{Lecture2\_Classifier\_Playground.html} & threat classification, thresholds, imagery detection (IoU)\\
3 & \texttt{Lecture3\_Regression\_Fusion\_Lab.html} & prediction intervals, Kalman tracking, casualty and force-ratio models\\
4 & \texttt{Lecture4\_Clustering\_Lab.html} & clustering intercepts, choosing $k$, DBSCAN, PCA\\
5 & \texttt{Lecture5\_Anomaly\_Detection\_Lab.html} & intrusion and sensor anomalies, alert budgets, drift detection\\
6 & \texttt{Lecture6\_Q\_Learning\_Gridworld.html} & patrol and routing agents, value iteration, DQN bias\\
7 & \texttt{Lecture7\_Policy\_Gradient\_Lab.html} & continuous control, PPO tuning, advantage estimation\\
8 & \texttt{Lecture8\_Multi\_Agent\_Lab.html} & team coordination, security games, credit, agent evaluation\\
9 & \texttt{Lecture9\_Revision\_Studio.html} & end-to-end decision chains and formula checks\\\bottomrule\end{tabularx}\end{center}
\begin{instr}{Suggested assessment rule}Credit a studio demonstration in the presentation only when the team states what they expected and why before revealing the result, and afterwards what the result implies for their system. Showing a studio's default numbers without changing an input adds nothing.\end{instr}

\needspace{12\baselineskip}\section{Quiz Administration and Answer Key}

\begin{sayit}``You will now take a 15-question quiz. Each question has four options and exactly one correct answer. You have 40 seconds per question; the quiz advances automatically when time runs out. A correct answer scores $+1$, a wrong answer $-0.25$, and an unanswered question 0. The maximum is 15 and the minimum is $-3.75$. Because of the negative marking, guessing blindly among four options gains nothing on average --- answer when you can eliminate options. Your score and a full review appear at the end.''\end{sayit}

\textbf{Delivery.} Open \texttt{Lecture10\_Quiz.html} on each student device (or project it and collect answers on paper). The page enforces the timer and computes the score with negative marking. Release \texttt{Lecture10\_Answers\_Assignment.html} only after all students have finished.

\subsection*{Answer key}\begin{center}\small\renewcommand{\arraystretch}{1.3}\begin{tabular}{|c|c|c|c|c|c|c|c|c|c|c|c|c|c|c|}\hline \textbf{1} & \textbf{2} & \textbf{3} & \textbf{4} & \textbf{5} & \textbf{6} & \textbf{7} & \textbf{8} & \textbf{9} & \textbf{10} & \textbf{11} & \textbf{12} & \textbf{13} & \textbf{14} & \textbf{15}\\\hline D & A & B & C & A & B & B & D & C & B & C & D & A & C & B\\\hline\end{tabular}\end{center}

\textbf{Interpreting scores.} 12 or more: secure command of the lecture. 9--11.75: good, review missed items. Below 9: revisit the reference text before the assignment. If more than half the class misses a question, revisit that point in the next lecture's opening Q\&A.

\subsection*{Questions, options and rationale}

{\small

\noindent\textbf{Q1.} Which problem statement is strongest?\begin{itemize}[leftmargin=2.2em,itemsep=0pt,topsep=1pt]\item[A.] We applied ML to air defence\item[B.] Our model is very accurate\item[C.] We used a deep network\item[\textbf{D.}] \textbf{IADS operators get 400 tracks/hour and need $>$=95\% accuracy with bomber recall $>$95\% in $<$=3 ms}~{\color{GrnN}$\checkmark$}\end{itemize}\noindent{\color{GreyN}\textit{Rationale:}} Specific, quantified, operational, with a clear decision supported.\par\medskip

\noindent\textbf{Q2.} A results slide must include at least:\begin{itemize}[leftmargin=2.2em,itemsep=0pt,topsep=1pt]\item[\textbf{A.}] \textbf{A baseline comparison, the appropriate per-class or task metrics and uncertainty}~{\color{GrnN}$\checkmark$}\item[B.] Only accuracy\item[C.] The training loss\item[D.] The code\end{itemize}\noindent{\color{GreyN}\textit{Rationale:}} 87\% means nothing without context.\par\medskip

\noindent\textbf{Q3.} Your accuracy is 95\% on 400 test cases. The approximate 95\% CI half-width is:\begin{itemize}[leftmargin=2.2em,itemsep=0pt,topsep=1pt]\item[A.] 0.2\%\item[\textbf{B.}] \textbf{2.1\%}~{\color{GrnN}$\checkmark$}\item[C.] 5\%\item[D.] 10\%\end{itemize}\noindent{\color{GreyN}\textit{Rationale:}} $1.96\sqrt{0.95\times0.05/400}=0.021$.\par\medskip

\noindent\textbf{Q4.} To compare two classifiers on the same test set you should use:\begin{itemize}[leftmargin=2.2em,itemsep=0pt,topsep=1pt]\item[A.] Two independent CIs only\item[B.] A t-test on training loss\item[\textbf{C.}] \textbf{McNemar's test on paired disagreements}~{\color{GrnN}$\checkmark$}\item[D.] Nothing\end{itemize}\noindent{\color{GreyN}\textit{Rationale:}} Paired predictions require a paired test.\par\medskip

\noindent\textbf{Q5.} For RL projects, results should be reported:\begin{itemize}[leftmargin=2.2em,itemsep=0pt,topsep=1pt]\item[\textbf{A.}] \textbf{As mean and spread over several seeds and opponents}~{\color{GrnN}$\checkmark$}\item[B.] From the single best seed\item[C.] As training reward only\item[D.] Without baselines\end{itemize}\noindent{\color{GreyN}\textit{Rationale:}} RL is high-variance; at least five seeds.\par\medskip

\noindent\textbf{Q6.} An ablation study:\begin{itemize}[leftmargin=2.2em,itemsep=0pt,topsep=1pt]\item[A.] Adds features\item[\textbf{B.}] \textbf{Removes one component at a time to show what drives performance}~{\color{GrnN}$\checkmark$}\item[C.] Tunes the learning rate\item[D.] Is optional decoration\end{itemize}\noindent{\color{GreyN}\textit{Rationale:}} It shows which parts of the system matter.\par\medskip

\noindent\textbf{Q7.} Which limitation statement is best?\begin{itemize}[leftmargin=2.2em,itemsep=0pt,topsep=1pt]\item[A.] Future work could improve accuracy\item[\textbf{B.}] \textbf{Trained only on L-band signatures; frequency-agile X-band radars would need retraining}~{\color{GrnN}$\checkmark$}\item[C.] There may be bugs\item[D.] More data would help\end{itemize}\noindent{\color{GreyN}\textit{Rationale:}} Specific and testable.\par\medskip

\noindent\textbf{Q8.} STARE stands for:\begin{itemize}[leftmargin=2.2em,itemsep=0pt,topsep=1pt]\item[A.] Search, Test, Analyse, Report, Evaluate\item[B.] Simulate, Tune, Apply, Retrain, Explain\item[C.] Start, Train, Assess, Run, Exit\item[\textbf{D.}] \textbf{State, Tie, Acknowledge, Reason, End}~{\color{GrnN}$\checkmark$}\end{itemize}\noindent{\color{GreyN}\textit{Rationale:}} A structured, honest way to answer hard questions.\par\medskip

\noindent\textbf{Q9.} When you do not know an answer, the panel prefers that you:\begin{itemize}[leftmargin=2.2em,itemsep=0pt,topsep=1pt]\item[A.] Bluff confidently\item[B.] Change the subject\item[\textbf{C.}] \textbf{Acknowledge the gap honestly and reason about what would happen}~{\color{GrnN}$\checkmark$}\item[D.] Cite an unrelated result\end{itemize}\noindent{\color{GreyN}\textit{Rationale:}} Honesty is rewarded; bluffing is detected.\par\medskip

\noindent\textbf{Q10.} Which rubric criterion carries the highest weight?\begin{itemize}[leftmargin=2.2em,itemsep=0pt,topsep=1pt]\item[A.] Problem definition\item[\textbf{B.}] \textbf{Algorithm justification and results/validation (25\% each)}~{\color{GrnN}$\checkmark$}\item[C.] Ethics\item[D.] Peer evaluation\end{itemize}\noindent{\color{GreyN}\textit{Rationale:}} 20/25/25/20/10.\par\medskip

\noindent\textbf{Q11.} Claiming 'ready for deployment' from simulation results is:\begin{itemize}[leftmargin=2.2em,itemsep=0pt,topsep=1pt]\item[A.] Acceptable\item[B.] Required\item[\textbf{C.}] \textbf{A common pitfall; the sim-to-real gap must be estimated and closed}~{\color{GrnN}$\checkmark$}\item[D.] Only wrong for RL\end{itemize}\noindent{\color{GreyN}\textit{Rationale:}} Drones fell 94\% to 78\% without randomisation.\par\medskip

\noindent\textbf{Q12.} Peer evaluation contributes:\begin{itemize}[leftmargin=2.2em,itemsep=0pt,topsep=1pt]\item[A.] 0\%\item[B.] 100\%\item[C.] 50\%\item[\textbf{D.}] \textbf{20\% of the project grade}~{\color{GrnN}$\checkmark$}\end{itemize}\noindent{\color{GreyN}\textit{Rationale:}} Moderated by the instructor.\par\medskip

\noindent\textbf{Q13.} A deployment plan should include:\begin{itemize}[leftmargin=2.2em,itemsep=0pt,topsep=1pt]\item[\textbf{A.}] \textbf{Autonomy level, latency, interface, ROE/IHL compliance, sim-to-real gap and monitoring}~{\color{GrnN}$\checkmark$}\item[B.] Only accuracy\item[C.] The dataset size only\item[D.] Marketing claims\end{itemize}\noindent{\color{GreyN}\textit{Rationale:}} This separates operationally mature projects from merely competent ones.\par\medskip

\noindent\textbf{Q14.} A trivial baseline for an imbalanced classifier is:\begin{itemize}[leftmargin=2.2em,itemsep=0pt,topsep=1pt]\item[A.] Random forest\item[B.] A deep network\item[\textbf{C.}] \textbf{Always predicting the majority class}~{\color{GrnN}$\checkmark$}\item[D.] SVM\end{itemize}\noindent{\color{GreyN}\textit{Rationale:}} Shows how much of the accuracy is just class frequency.\par\medskip

\noindent\textbf{Q15.} Why split sequential data by mission or day rather than randomly?\begin{itemize}[leftmargin=2.2em,itemsep=0pt,topsep=1pt]\item[A.] To save time\item[\textbf{B.}] \textbf{To prevent temporal leakage between correlated samples in train and test}~{\color{GrnN}$\checkmark$}\item[C.] It is required by law\item[D.] To balance classes\end{itemize}\noindent{\color{GreyN}\textit{Rationale:}} Random splits put near-duplicates on both sides.\par\medskip

}

\needspace{12\baselineskip}\section{Assignment: Marking Scheme and Model Solutions}

\textbf{Instructions.} This assignment rehearses your project. 100 marks. Model answers are provided for self-assessment.

\begin{center}\small\renewcommand{\arraystretch}{1.25}\begin{tabularx}{0.85\textwidth}{@{}l X r@{}}\toprule\textbf{\#} & \textbf{Problem} & \textbf{Marks}\\\midrule1 & Project proposal & 40\\
2 & Results with uncertainty & 30\\
3 & Rehearsing the panel & 30\\\midrule & \textbf{Total} & \textbf{100}\\\bottomrule\end{tabularx}\end{center}

\begin{instr}{General marking guidance}\begin{itemize}
\item Marks for each part are shown in square brackets in the problem statement.
\item Award method marks for correct working even when an arithmetic slip carries through; deduct once for the slip, not again for its consequences.
\item For interpretive parts, full marks need (i) the correct conclusion, (ii) the reason, and (iii) an operational implication or caveat.
\item Answers that state a number without uncertainty, baseline or caveat where one is required lose a mark per part.
\item Suggested deadline: one week after the lecture; return marked work before the next quiz on the same topic.
\end{itemize}\end{instr}

\needspace{8\baselineskip}\subsection{Problem 1: Project proposal \normalfont\small\color{GreyN}(40 marks)}

\begin{prob}Choose one catalogue topic (T1--T8) and write a one-page proposal: problem statement with numbers; data and its pitfalls; candidate methods with at least two alternatives; evaluation protocol with baselines and uncertainty; deployment concept; two specific limitations.\end{prob}

\begin{soln}\textbf{Model answer (T4, cyber anomaly):} \textit{Problem:} a brigade network generates 2M events/day; analysts can triage 200 alerts; the signature IDS misses novel attacks. \textit{Data:} 30 days of normal traffic plus red-team attacks; split by day; base rate about 0.01\%; possible survivorship (only detected incidents labelled) and drift. \textit{Methods:} Isolation Forest screen + autoencoder confirmation; alternatives: z-score/Mahalanobis, one-class SVM. \textit{Evaluation:} TPR at a fixed 150 false alerts/day; precision at the base rate; bootstrap CIs; per-attack recall; ablation without tier 1; a held-out novel attack type. \textit{Deployment:} edge server, Level-2 containment approvals, drift monitoring, monthly retraining. \textit{Limitations:} mimicry attacks imitating normal traffic; red-team traffic may not represent real adversaries.\end{soln}

\needspace{8\baselineskip}\subsection{Problem 2: Results with uncertainty \normalfont\small\color{GreyN}(30 marks)}

\begin{prob}Your classifier scores 93.0\% on 500 test tracks; a baseline scores 91.0\% on the same tracks. Of the tracks on which they disagree, yours is right on 34 and the baseline on 24.

(a) Compute 95\% CIs for both accuracies. [10]\newline (b) Apply McNemar's test. [10]\newline (c) Write the results sentence you would put on your slide. [10]\end{prob}

\begin{soln}\textbf{(a)} Yours: $1.96\sqrt{0.93(0.07)/500}=2.2\%$ $\rightarrow$ (90.8, 95.2). Baseline: $1.96\sqrt{0.91(0.09)/500}=2.5\%$ $\rightarrow$ (88.5, 93.5). The intervals overlap.

\textbf{(b)} $b=34$, $c=24$: $\chi^2=(34-24)^2/58=1.72$; $p\approx0.19$. Not significant.

\textbf{(c)} ``Our classifier achieved 93.0\% (95\% CI 90.8--95.2) versus 91.0\% for the baseline on the same 500 tracks; the difference is not statistically significant (McNemar $p=0.19$), so a larger test set or deployment-domain data is needed to establish an advantage.''\end{soln}

\needspace{8\baselineskip}\subsection{Problem 3: Rehearsing the panel \normalfont\small\color{GreyN}(30 marks)}

\begin{prob}Write STARE answers to three likely panel questions for your project: one on data bias, one on false-alarm cost, and one on autonomy and accountability.\end{prob}

\begin{soln}\textbf{Data bias.} S: trained on 30 days of traffic from two segments. T: slide 3 lists the sources. A: no data from deployed forward networks. R: those networks may show different protocols, so the FPR may rise. E: collect a 2-week sample from the target network and re-calibrate the threshold.

\textbf{False-alarm cost.} S: 150 false alerts/day at 87\% TPR. T: slide 6 trade-off curve. A: analyst time per alert is assumed, not measured. R: if triage takes longer than 2 min, capacity is exceeded. E: measure triage time and add per-host aggregation.

\textbf{Autonomy.} S: the system recommends isolation of hosts. T: slide 8 interface mock-up. A: automatic isolation was not tested for collateral impact. R: false isolation could cut critical services. E: keep Level 2 (analyst approval), with full logging and a rollback procedure.\end{soln}

\end{document}